% =====================================================================
%  14  全体最適化 — 左=手法／評価関数(TAC・CAPEX・OPEX の定義と分類)／制約、
%   右=設計変数22の探索範囲＋最適化で決定した値(赤)。F_freshは従属変数のため非掲載。
%  source_memo: 評価関数・CAPEX/OPEX 定義=03_economic_indicators.tex（TAC=CAPEX/n+OPEX, n=8、
%   OPEX=0.180 C_TM+2.73 C_OL+1.23(C_UT+C_WT+C_RM)）、探索範囲=Z:\pdh_simulator\main.py SEARCH_SPACE、
%   決定値=02_process_overview.tex Table 2.3 / best.json trial #201（スライド3と一致）、
%   制約=Z:\pdh_simulator\flowsheet\specs.py＋units\reactors\catofin.py。
%  方針：帯と本文を重ねない（上 vspace は小さな正）。余白を先に使い、足りない分だけ縮小。
%        評価関数を主役に TAC/CAPEX/OPEX の定義・分類を詳述。決定値は赤＝最適化の出力。丸括弧禁止。
% =====================================================================
\definecolor{SpGreen}{HTML}{1FB24D}   % 製品目標スペック強調用の明るい緑
\begin{frame}[t]

  \vspace*{0.2ex}
  \begin{columns}[T,onlytextwidth]
    % ===== 左：最適化手法・評価関数(TAC/CAPEX/OPEX 詳述)・制約条件 ====
    \begin{column}{0.53\textwidth}
      \fbox{\scriptsize 最適化手法}\;{\footnotesize\bfseries ベイズ最適化(TPE)}\par

      \vspace{0.34em}
      \fbox{\scriptsize 評価関数}\;{\footnotesize 年間総費用 {\color{SpRed}\bfseries TAC} を最小化}\par
      \vspace{0.2em}
      {\normalsize\boldmath$ {\color{SpRed}\mathrm{TAC}}=\mathrm{CAPEX}/n+\mathrm{OPEX} $\unboldmath}\par

      \vspace{0.46em}
      {\raggedright
       {\footnotesize\bfseries CAPEX}\par
       {\scriptsize\linespread{1.05}\selectfont 資本的支出　建設の初期投資\par
        \hspace{0.9em}直接費：機器費・据付・配管・計装・土木\par
        \hspace{0.9em}間接費：設計・予備費・エンジニアリング\par}
       \vspace{0.38em}
       {\footnotesize\bfseries OPEX}\par
       {\scriptsize\linespread{1.05}\selectfont 事業運営費　年間の運転費\par
        \hspace{0.9em}変動費：原料費・用役費・廃棄物処理費\par
        \hspace{0.9em}固定費：労務費・保全費・一般経費\par
        \vspace{0.95em}
        $n$：法定耐用 $8$ 年\par}
      }

      \vspace{0.46em}
      \fbox{\scriptsize 制約条件}\par
      \vspace{0.18em}
      {\scriptsize\linespread{1.0}\selectfont
       \begin{tabular}{@{}p{0.47\linewidth}@{\hspace{0.3em}}|@{\hspace{0.3em}}p{0.50\linewidth}@{}}
         \raggedright
         {\color{SpGreen}・C3H6 純度 $\ge 99.5$ wt\%}\par
         ・H2 純度 $\ge 99.9$ mol\%\par
         {\color{SpGreen}・生産量 $\ge 400$ kt/年}\par
         ・反応器 床$\Delta P\le 5$\%\par
         ・反応器 総$\Delta P\le 10$\%\par
         ・空塔速度 $\ge 0.15$ m/s
         &
         \raggedright
         ・触媒体積 $\le 200\ \mathrm{m^3}$／基\par
         ・PSA 床圧損 $\le 0.3$ bar\par
         ・脱エタン塔頂が凝縮可\par
         ・膜供給圧 $>$ フィード圧\par
         ・リサイクル収支整合
       \end{tabular}}
    \end{column}

    \begin{column}{0.01\textwidth}\end{column}   % 中央の余白

    % ===== 右：設計変数22の探索範囲＋決定値(赤)。ユニットごとに枠を分割 ======
    \begin{column}{0.46\textwidth}
      \setlength{\fboxsep}{1.5pt}\setlength{\fboxrule}{0.8pt}%
      \fcolorbox{HeaderBlue}{white}{\makebox[\dimexpr\linewidth-2\fboxsep-2\fboxrule\relax][l]{%
        \scriptsize\renewcommand{\arraystretch}{0.82}%
        \begin{tabular}{@{}l@{\hspace{0.5em}}l@{\hspace{0.5em}}l@{\hspace{0.6em}}r@{\hspace{0.35em}}l@{}}
          入口温度     & $T_{\mathrm{in}}$  & $930$–$940$   & {\color{SpRed}$935.2$} & K \\
          サイクル時間 & $t_{\mathrm{cyc}}$ & $12$–$25$     & {\color{SpRed}$14.0$}  & min \\
          反応器内径   & $D$                & $9$–$12$      & {\color{SpRed}$10.76$} & m \\
          浅床厚       & $L_{\mathrm{bed}}$ & $0.3$–$3.0$   & {\color{SpRed}$1.58$}  & m \\
          並列基数     & $N_{\mathrm{on}}$  & $6$–$8$       & {\color{SpRed}$7$}     & --- \\
          触媒粒径     & $d_p$              & $4$–$6$       & {\color{SpRed}$5.84$}  & mm \\
        \end{tabular}}}\par\vspace{0.6pt}
      \fcolorbox{HeaderBlue}{white}{\makebox[\dimexpr\linewidth-2\fboxsep-2\fboxrule\relax][l]{%
        \scriptsize\renewcommand{\arraystretch}{0.82}%
        \begin{tabular}{@{}l@{\hspace{0.5em}}l@{\hspace{0.5em}}l@{\hspace{0.6em}}r@{\hspace{0.35em}}l@{}}
          PSA 塔径     & $D_{\mathrm{P}}$   & $2.9$–$5.0$   & {\color{SpRed}$4.49$}  & m \\
          PSA 層高     & $L_{\mathrm{P}}$   & $22$–$30$     & {\color{SpRed}$24.8$}  & m \\
          脱着基準     & $\beta$            & $0.22$–$0.40$ & {\color{SpRed}$0.254$} & --- \\
        \end{tabular}}}\par\vspace{0.6pt}
      \fcolorbox{HeaderBlue}{white}{\makebox[\dimexpr\linewidth-2\fboxsep-2\fboxrule\relax][l]{%
        \scriptsize\renewcommand{\arraystretch}{0.82}%
        \begin{tabular}{@{}l@{\hspace{0.5em}}l@{\hspace{0.5em}}l@{\hspace{0.6em}}r@{\hspace{0.35em}}l@{}}
          膜供給圧     & $P_H$              & $7.5$–$9.5$   & {\color{SpRed}$8.44$}  & bar \\
          膜面積       & $A_{\mathrm{mem}}$ & $0.5$–$3$     & {\color{SpRed}$1.28$}  & $\times10^{5}\,\mathrm{m^2}$ \\
        \end{tabular}}}\par\vspace{0.6pt}
      % 原料流量 F_fresh は最適化変数ではなく従属変数（製品スペック固定下で物質収支から決定、
      % 決定値 1494.6 kmol/h）のため本表から除外（22 変数）。
      \fcolorbox{HeaderBlue}{white}{\makebox[\dimexpr\linewidth-2\fboxsep-2\fboxrule\relax][l]{%
        \scriptsize\renewcommand{\arraystretch}{0.82}%
        \begin{tabular}{@{}l@{\hspace{0.5em}}l@{\hspace{0.5em}}l@{\hspace{0.6em}}r@{\hspace{0.35em}}l@{}}
          塔圧力       & $P_1$ & $16$–$20$    & {\color{SpRed}$19.5$} & bar \\
                       & $P_2$ & $7.7$–$8.2$  & {\color{SpRed}$7.88$} & bar \\
                       & $P_3$ & $16$–$19$    & {\color{SpRed}$16.8$} & bar \\
          塔段数       & $N_1$ & $30$–$60$    & {\color{SpRed}$36$}   & --- \\
                       & $N_2$ & $60$–$80$    & {\color{SpRed}$61$}   & --- \\
                       & $N_3$ & $115$–$160$  & {\color{SpRed}$117$}  & --- \\
          フィード段   & $N_{f1}$ & $22$–$28$ & {\color{SpRed}$28$}   & --- \\
          フィード比   & $f_2$ & $0.40$–$0.60$ & {\color{SpRed}$0.51$} & --- \\
                       & $f_3$ & $0.60$–$0.90$ & {\color{SpRed}$0.82$} & --- \\
          組成分率     & $\phi_1$ & $0.90$–$0.999$ & {\color{SpRed}$0.995$} & --- \\
          還流比       & $R_2$ & $10$–$15$    & {\color{SpRed}$11.8$} & --- \\
        \end{tabular}}}\par
      \vspace{1pt}
      {\tiny {\color{SpRed}\bfseries 赤：決定値}・$1$脱ブタン $2$脱エタン $3$C3塔}
    \end{column}
  \end{columns}

\end{frame}
